\par\Needspace{6\baselineskip}\medskip
\hypertarget{native-159-article114}{}\pdfbookmark[2]{114.}{article-114}
\begin{center}114.\end{center}
By multiplying equation 1 by $\sin\alpha''\operatorname{tang}B''-\sin L''\operatorname{tang}\beta''$, equation 2 by $\cos L''\operatorname{tang}\beta''-\cos\alpha''\operatorname{tang}B''$, and equation 3 by $\sin(L''-\alpha'')$, and adding the products, we obtain
\[
\begin{aligned}
0={}&n\{(0.2.\mathrm{II})\delta+(\mathrm{O}.2.\mathrm{II})D\}\\
&-n'\{(1.2.\mathrm{II})\delta'+(\mathrm{I}.2.\mathrm{II})D'\}
\end{aligned}\tag*{4)}
\]
and similarly, or more conveniently by simply permuting the positions,
\[
\begin{aligned}
0={}&n\{(0.1.\mathrm{I})\delta+(\mathrm{O}.1.\mathrm{I})D\}\\
&+n''\{(2.1.\mathrm{I})\delta''+(\mathrm{II}.1.\mathrm{I})D''\}
\end{aligned}\tag*{5)}
\]
\[
\begin{aligned}
0={}&n'\{(1.0.\mathrm{O})\delta'+(\mathrm{I}.0.\mathrm{O})D'\}\\
&-n''\{(2.0.\mathrm{O})\delta''+(\mathrm{II}.0.\mathrm{O})D''\}.
\end{aligned}\tag*{6)}
\]
Thus, if the ratio of the quantities $n$, $n'$ is given, equation 4 lets us use $\delta$ to determine $\delta'$, or obtain $\delta$ from $\delta'$. The same holds for equations 5 and 6. Combining equations 4, 5, and 6 gives the equation
\[
\begin{aligned}
&\frac{(0.2.\mathrm{II})\delta+(\mathrm{O}.2.\mathrm{II})D}{(0.1.\mathrm{I})\delta+(\mathrm{O}.1.\mathrm{I})D}\\[.6em]
&\qquad\times\frac{(1.0.\mathrm{O})\delta'+(\mathrm{I}.0.\mathrm{O})D'}{(1.2.\mathrm{II})\delta'+(\mathrm{I}.2.\mathrm{II})D'}\\[.6em]
&\qquad\times\frac{(2.1.\mathrm{I})\delta''+(\mathrm{II}.1.\mathrm{I})D''}{(2.0.\mathrm{O})\delta''+(\mathrm{II}.0.\mathrm{O})D''}=-1,
\end{aligned}\tag*{7)}
\]
which allows us to determine the third distance of the celestial body from the Earth from two such distances. It can be shown, however, that equation 7 becomes an identity and is therefore unsuitable for determining one distance from the other two whenever $B=B'=B''=0$ and
\[
\left.\begin{aligned}
&\operatorname{tang}\beta'\operatorname{tang}\beta''\sin(L-\alpha)\sin(L''-L')\\
&+\operatorname{tang}\beta''\operatorname{tang}\beta\sin(L'-\alpha')\sin(L-L'')\\
&+\operatorname{tang}\beta\operatorname{tang}\beta'\sin(L''-\alpha'')\sin(L'-L)
\end{aligned}\right\}=0.
\]
\par
The following formula is free of this difficulty and is \SourceLeaf{160}readily derived from equations 1, 2, and 3.
\[
\begin{aligned}
&(0.1.2)\delta\delta'\delta''+(\mathrm{O}.1.2)D\delta'\delta''\\
&+(0.\mathrm{I}.2)D'\delta\delta''+(0.1.\mathrm{II})D''\delta\delta'\\
&+(0.\mathrm{I}.\mathrm{II})D'D''\delta+(\mathrm{O}.1.\mathrm{II})DD''\delta'\\
&+(\mathrm{O}.\mathrm{I}.2)DD'\delta''+(\mathrm{O}.\mathrm{I}.\mathrm{II})DD'D''=0.
\end{aligned}\tag*{8)}
\]
\par
By multiplying equation 1 by $\sin\alpha'\operatorname{tang}\beta''-\sin\alpha''\operatorname{tang}\beta'$, equation 2 by $\cos\alpha''\operatorname{tang}\beta'-\cos\alpha'\operatorname{tang}\beta''$, and equation 3 by $\sin(\alpha''-\alpha')$, and adding the products, we obtain
\[
\begin{aligned}
0={}&n\{(0.1.2)\delta+(\mathrm{O}.1.2)D\}\\
&-n'(\mathrm{I}.1.2)D'+n''(\mathrm{II}.1.2)D''
\end{aligned}\tag*{9)}
\]
and likewise
\[
\begin{aligned}
0={}&n(0.\mathrm{O}.2)D-n'\{(0.1.2)\delta'+(0.\mathrm{I}.2)D'\}\\
&+n''(0.\mathrm{II}.2)D''
\end{aligned}\tag*{10)}
\]
\[
\begin{aligned}
0={}&n(0.1.\mathrm{O})D-n'(0.1.\mathrm{I})D'\\
&+n''\{(0.1.2)\delta''+(0.1.\mathrm{II})D''\}.
\end{aligned}\tag*{11)}
\]
Given the ratio among the quantities $n$, $n'$, $n''$, we can use these equations to derive the distances $\delta$, $\delta'$, $\delta''$. This conclusion, however, holds only in general and admits an exception whenever $(0.1.2)=0$. In this case, it can be shown that equations 9, 10, and 11 yield only a necessary relation among the quantities $n$, $n'$, $n''$, and that each of the three yields the same relation. The analogous restrictions on equations 4, 5, and 6 will readily occur to the experienced reader.
\par
Finally, all the conclusions developed here are useless whenever the orbital plane coincides with the ecliptic. For if $\beta$, $\beta'$, $\beta''$, $B$, $B'$, $B''$ are all $=0$, equation 3 is an \emph{identity}, and hence so are all the equations that follow it.
\par\bigskip\begin{center}\rule{.23\linewidth}{.35pt}\end{center}
